\begin{center}
    \large
    \begin{singlespace}    
        \textbf{\chineseTitle{}} \\[0.5cm]
    \end{singlespace}

    \begin{singlespace}    
    學生：\studentCnName  \hspace{2.5cm}  指導教授：\advisorCnName \hspace{0.1cm} 博士 \\

    \ifdefined\advisorCnNameB % 如果有共同指導教授
        \hspace{9.6cm}  \advisorCnNameB ~ 博士 \\
    \fi
    \end{singlespace}

    \vspace{0.5cm}

    \begin{singlespace}
        國立陽明交通大學\NameofDepartmentInstituteCN\\[0.2cm]
    \end{singlespace}

    \textbf{摘\hspace{1cm}要} \\[0.5cm]

\end{center}

\normalsize

\begin{spacing}{1.45}
知識圖譜感知推薦在稠密且品質良好的圖譜上通常能提升準確度；然而，在評論衍生、冷啟動或隱私受限等情境中，圖譜往往稀疏或品質不穩，相關研究仍相對不足。本文以商品評論衍生的稀疏圖譜為例，過濾後每件商品平均僅有 2.4 條面向邊；實驗顯示，代表性的知識圖譜感知方法皆不如完全不使用圖譜的純協同過濾模型。主要原因在於，這些方法將圖譜訊號硬性嵌入評分流程，缺乏在圖譜不可靠時自動降低其影響的機制，使雜訊干擾協同訊號的學習。

本文提出 RA-GARK，將圖譜視為一條可調控的側通道，而非評分流程中必須使用的組件，並透過三個步驟實現。其一，將每件商品的圖譜內容整理為少數代表性面向，並以帶有語意的向量初始化，避免完全隨機的起點。其二，模型會依使用者與商品，自動選取最相關的面向。其三，由融合閘決定是否啟用此側通道：訓練初期幾乎關閉，使模型從純協同過濾出發；只有當側通道確實有幫助時才逐步打開，並在其無用時平順退回。在不更動圖譜本身的前提下，RA-GARK 將圖譜的貢獻由負面轉為正面。

本研究指出，知識圖譜感知方法在稀疏圖譜上表現不佳，關鍵並非圖譜完全沒有資訊，而是既有架構缺乏選擇性使用圖譜的能力。同一個圖譜在強制融合下可能成為負擔，但當模型能自行決定是否採用時，便可轉為助益。這種優雅退化的設計並不限於知識圖譜；對於任何需要結合可靠主訊號與品質不穩定輔助訊號的情境，例如冷啟動、多模態或跨領域推薦，同樣具有應用價值。\\[0.5\baselineskip]
\textbf{關鍵字：知識圖譜推薦；可退化機制；融合閘；稀疏知識圖譜；協同過濾}
\end{spacing}
